\section{Исполнимая маршрутизация и аудит разработки}
Структурное исследование использует новый зафиксированный протокол и новые случайные потоки обучения, валидации, планирования и тестирования. Прежняя независимая реализация сохранена в версионированном исследовательском архиве. Её шесть отдельных сравнений при routing feedback и известной структуре не установили положительного эффекта после поправки на множественные проверки; описательная нормированная разность с равными весами процессов составила $-0.090$ при заданной структуре и $0.001$ при обнаруженной. Эти исходы не объединяются с новыми основными сравнениями.

До новой фиксации аналитический аудит выявил две ошибки исполнителя. Два дерева с наблюдаемым языком $(BD)^+$ размещали решение о возврате в разных точках. При фиксированной вероятности возврата $0.2$ прежний исполнитель давал 1.250 и 2.027 прохода на экземпляр; уменьшение вероятности вдвое давало 1.1085 и неизменные 2.027. Проверка содержит 2000 независимо сгенерированных экземпляров. Преобразование общего продолжения переносит выбор в наблюдаемую точку решения и обеспечивает точное совпадение двух массивов траекторий с вмешательством и без него. Ожидаемые числа проходов равны $1/0.8$ и $1/0.9$. Отдельная проверка XOR с условием выявила 999 нарушений среди 2000 экземпляров: имитация обученной модели игнорировала явные условия по классу. Исправленный исполнитель не допускает нарушений. Эти проверки устанавливают свойства исполнения, а не величину эффекта по стоимости политики.

Отдельный аудит из 92 запусков охватывает каждый процесс, механизм обратной связи и тип вмешательства. Отсчёты истинного генератора, записи жизненного цикла, завершения и времена активации идентичны до и после исправления. Уравнения генерации, мощности, стоимости, условия поступлений и масштабы нагрузки не меняются. Технические пилоты проверяют все шесть обнаруженных структур и эквивалентность планировочных оценок при известной и обнаруженной структуре; они исключены из статистических выводов.

Процедура обнаружения сравнивает стандартный Inductive Miner с кандидатом, который отдаёт приоритет существующему правилу строгого цикла и выводит детерминированные условия по наблюдаемым атрибутам. Допускаются только кандидаты, которым соответствуют все обучающие трассы по оптимальным выравниваниям и которые проходят проверки компилятора; затем максимизируется обучающая precision по token replay. Глобальная alignment fitness может быть чуть ниже единицы из-за ненулевой стоимости невидимых переходов; допуск использует долю соответствующих трасс. Кандидат не получает исходное дерево. Эквивалентность наблюдаемых языков проверяется только затем как диагностика на синтетических данных. Сохранены все кандидатные модели, условия, результаты проверки соответствия и квитанции подготовки.

Исполненная реализация ограничивает модель 30 наблюдаемыми метками, 62 позициями, 120 переходами и 32 ресурсными пулами при бинарных наблюдаемых атрибутах класса и процесса. Неподдерживаемые развилки и небезопасные маркировки отклоняются. Это ограничения реализации, а не эмпирическое утверждение о масштабируемости.

Дополнительная техническая проверка использует только исключённые из оценки пилоты. Общий объект конфигурации содержит поля истинного генератора, хотя подготовка на обученной модели их не использует. Подмена истинного дерева и генеративных коэффициентов, а также удаление топологии из входов обучения поведения оставляют обученные параметры, выбранную обнаруженную структуру и исполнение обученной модели полностью неизменными во всех шести пилотах. Ресурсные отображения, потенциальные точки возврата и наблюдаемые очереди остаются инженерными входами.
